\subsection{Probes Output}%
\label{SEC:Files:Probes_Output}
\key*[key:probes]{}[probes]

%%%%%%%%%%%%%%%%%%%%%%%%%%%%%%%%%%%%%%%%%%%%%%%%%%%%%%%%%%%%%%%%%%%%%%%%%%%%%%%%

With time accurate simulations, it is impractical to store the entire unsteady
solution history for every grid point.
%
This is why the GFR code primarily outputs instantaneous and/or time-averaged
files at either fixed iteration or time intervals.
%
However, it is sometimes convenient to have the entire unsteady flow history at
a limited subset of points for later analysis.
%
The GFR code provides the ability to output this unsteady history using point
probes.
%
Probes can either be specified as an individual probe at a specific coordinate
point using the input \inp[nmlOutput::PointProbes]{Point\_Probes}, as a rake of
probes that are evenly spaced along a line between two coordinate points using
the input \inp[nmlOutput::RakeProbes]{Rake\_Probes}, or the input
\inp[nmlOutpu::CellProbes]{Cell\_Probes} can be used to create a set of probes
consisting of every solution point within the grid cell containing the specified
coordinate point.
%
For cell probes, individual point probes will be created for every solution
point within each identified grid cell.
%
Note that for the point and rake probes, the solution point closest to each
requested probe location will be used for the source of the probe data; this is
done to reduce the performance overhead that otherwise would be incurred from
having to interpolate the solution to the requested probe location if it was not
located exactly at a solution point.
%
If a solution point is identified as the source of a probe by more than one
point\slash rake\slash cell probe, the duplicate probes will be removed and only
a single probe will be collected for that solution point.

%%%%%%%%%%%%%%%%%%%%%%%%%%%%%%%%%%%%%%%%%%%%%%%%%%%%%%%%%%%%%%%%%%%%%%%%%%%%%%%%

Each point probe is written to a separate Tecplot ASCII formatted file within
the directory \file{Probe\_files}.
%
Time-history files for probes specified as individual point probes will be named
\file{point-probe.$i$.tec}, where $i$ corresponds to the index of
\inp[nmlOutput::PointProbes]{Point\_Probes} from the input file.
%
Time-history files for probes within a rake of probes will be named
\file{rake-probe.$i$-$j$.tec}, where $i$ corresponds to the index of
\inp[nmlOutput::RakeProbes]{Rake\_Probes} from the input file, and $j$
corresponds to the index of the individual probe within rake $i$.
%
Time-history files for probes from a set of cell probes will be named
\file{cell-probe.$i$-$j$.tec}, where $i$ corresponds to the index of the grid
cell within the global grid, and $j$ corresponds to the solution point index of
the individual probe which will be between 1 and the total number of solution
points within cell $i$.
%
Each file contains a simple header and comments to identify the requested probe
location from the input, the actual probe location, and flow variables.

%%%%%%%%%%%%%%%%%%%%%%%%%%%%%%%%%%%%%%%%%%%%%%%%%%%%%%%%%%%%%%%%%%%%%%%%%%%%%%%%

The data for all probes is sampled at an iteration interval that is specified by
the input
\inp[nmlOutput::ProbesCollectionInterval]{Probes\_Collection\_Interval} and
saved to a buffer.
%
The input variable \inp[nmlOutput::ProbesWriteInterval]{Probes\_Write\_Interval}
specifies how many data collections to perform until writing out the buffer to
the individual probe time-history files.

%%%%%%%%%%%%%%%%%%%%%%%%%%%%%%%%%%%%%%%%%%%%%%%%%%%%%%%%%%%%%%%%%%%%%%%%%%%%%%%%

The time-history probe files are appended each time the code is executed.
%
If the number or order of probes changes between two consecutive simulations,
the user should remove the existing files before proceeding.

%%%%%%%%%%%%%%%%%%%%%%%%%%%%%%%%%%%%%%%%%%%%%%%%%%%%%%%%%%%%%%%%%%%%%%%%%%%%%%%%

If time-history data at an evenly spaced time interval is required, then a
constant time step should be used.

%%%%%%%%%%%%%%%%%%%%%%%%%%%%%%%%%%%%%%%%%%%%%%%%%%%%%%%%%%%%%%%%%%%%%%%%%%%%%%%%

\SeeAlso{
\inp[nmlOutput::ProbesCollectionInterval]{Probes\_Collection\_Interval},
\inp[nmlOutput::ProbesWriteInterval]{Probes\_Write\_Interval},
\inp[nmlOutput::PointProbes]{Point\_Probes},
\inp[nmlOutput::RakeProbes]{Rake\_Probes},
\inp[nmlOutput::RakeProbes]{Cell\_Probes}
}

